\documentclass[10pt,a4paper]{article}

\usepackage[utf8]{inputenc}
\usepackage[T1]{fontenc}
\usepackage[margin=1.5cm]{geometry}
\usepackage{amsmath,amssymb}
\usepackage{tikz}
\usetikzlibrary{positioning,calc,arrows.meta,matrix}
\usepackage{xcolor}
\usepackage{listings}
\usepackage{tcolorbox}
\tcbuselibrary{breakable}
\usepackage{enumitem}

\lstset{
    language=C,
    basicstyle=\ttfamily\scriptsize,
    backgroundcolor=\color{gray!10},
    keywordstyle=\color{blue},
    commentstyle=\color{gray},
    stringstyle=\color{red},
    numbers=left,
    numberstyle=\tiny,
    breaklines=true,
    frame=single,
    tabsize=4,
    showstringspaces=false,
    inputencoding=utf8,
    extendedchars=true,
    literate={é}{{\'e}}1 {è}{{\`e}}1 {à}{{\`a}}1 {ç}{{\c{c}}}1 
             {ê}{{\^e}}1 {î}{{\^i}}1 {ô}{{\^o}}1 {û}{{\^u}}1 {ù}{{\`u}}1
}

\newtcolorbox{exercicebox}[1]{colback=green!5,colframe=green!50!black,fonttitle=\bfseries,title=#1,breakable}

\title{\textbf{TD Chapitre 1 : Représentation Numérique}\\
\large Conversions, Entiers et Flottants}
\author{Dr. El Hadji Bassirou TOURÉ\\
Département de Mathématiques et Informatique\\
Université Cheikh Anta Diop de Dakar}
\date{}

\begin{document}
\maketitle

%==============================================================================
\section{Conversions de Base}
%==============================================================================

\begin{exercicebox}{Exercice 1 : Conversions binaire $\leftrightarrow$ hexadécimal $\leftrightarrow$ décimal}

Compléter le tableau suivant (valeurs non signées sur 8 bits) :

\begin{center}
\begin{tabular}{|c|c|c|}
\hline
\textbf{Binaire (8 bits)} & \textbf{Hexadécimal} & \textbf{Décimal} \\
\hline
$10110101$ & & \\
\hline
& \texttt{0x7F} & \\
\hline
& & 200 \\
\hline
$11001010$ & & \\
\hline
& \texttt{0xE3} & \\
\hline
& & 47 \\
\hline
\end{tabular}
\end{center}
\end{exercicebox}

\begin{exercicebox}{Exercice 2 : Conversion décimal $\rightarrow$ binaire par divisions successives}

En utilisant la méthode des divisions successives par 2, convertir en binaire sur 8 bits :

\begin{enumerate}[label=\alph*)]
    \item $73_{10}$
    \item $156_{10}$
    \item $241_{10}$
\end{enumerate}

Pour chaque conversion, montrer toutes les étapes (divisions et restes).
\end{exercicebox}

\begin{exercicebox}{Exercice 3 : Conversions rapides hex $\leftrightarrow$ binaire}

Convertir directement (sans passer par le décimal) :

\begin{enumerate}[label=\alph*)]
    \item \texttt{0xAF} en binaire
    \item \texttt{0x3C} en binaire
    \item $11110001_2$ en hexadécimal
    \item $01011110_2$ en hexadécimal
\end{enumerate}

\textit{Rappel : chaque chiffre hex correspond exactement à 4 bits.}
\end{exercicebox}

%==============================================================================
\section{Entiers Signés et Non Signés}
%==============================================================================

\begin{exercicebox}{Exercice 4 : Formule B2U (Binary to Unsigned)}

En utilisant la formule $B2U_w(\vec{x}) = \sum_{i=0}^{w-1} x_i \cdot 2^i$, calculer la valeur décimale des motifs binaires suivants interprétés comme \textbf{unsigned} sur 8 bits :

\begin{enumerate}[label=\alph*)]
    \item $10000001_2$
    \item $11111111_2$
    \item $01010101_2$
\end{enumerate}
\end{exercicebox}

\begin{exercicebox}{Exercice 5 : Formule B2T (Binary to Two's complement)}

En utilisant la formule $B2T_w(\vec{x}) = -x_{w-1} \cdot 2^{w-1} + \sum_{i=0}^{w-2} x_i \cdot 2^i$, calculer la valeur décimale des mêmes motifs interprétés comme \textbf{signed} sur 8 bits :

\begin{enumerate}[label=\alph*)]
    \item $10000001_2$
    \item $11111111_2$
    \item $01010101_2$
\end{enumerate}
\end{exercicebox}

\begin{exercicebox}{Exercice 6 : Négation en complément à deux}

Calculer la négation des nombres suivants sur 8 bits en utilisant $-x = \sim x + 1$ :

\begin{enumerate}[label=\alph*)]
    \item $x = 42$ (écrire d'abord en binaire)
    \item $x = 1$
    \item $x = -1$ (soit $11111111_2$)
    \item $x = -128$ (soit $10000000_2$) — que remarquez-vous ?
\end{enumerate}

Pour chaque cas, montrer : $x$ en binaire, $\sim x$, puis $\sim x + 1$.
\end{exercicebox}

\begin{exercicebox}{Exercice 7 : Plages de valeurs}

\begin{enumerate}[label=\alph*)]
    \item Donner $T_{min}$ et $T_{max}$ pour un entier signé sur 8 bits.
    \item Donner $U_{min}$ et $U_{max}$ pour un entier non signé sur 8 bits.
    \item Donner $T_{min}$ et $T_{max}$ pour \texttt{int16\_t} (16 bits signés).
    \item Pourquoi y a-t-il une asymétrie ($|T_{min}| \neq |T_{max}|$) pour les signés ?
\end{enumerate}
\end{exercicebox}

%==============================================================================
\section{Conversions Signed $\leftrightarrow$ Unsigned}
%==============================================================================

\begin{exercicebox}{Exercice 8 : Même motif, deux interprétations}

Pour chaque motif binaire sur 8 bits, donner sa valeur en unsigned ET en signed :

\begin{center}
\begin{tabular}{|c|c|c|}
\hline
\textbf{Binaire} & \textbf{Unsigned} & \textbf{Signed} \\
\hline
$01111111$ & & \\
\hline
$10000000$ & & \\
\hline
$11111111$ & & \\
\hline
$11111110$ & & \\
\hline
$11000000$ & & \\
\hline
\end{tabular}
\end{center}
\end{exercicebox}

\begin{exercicebox}{Exercice 9 : Formules de conversion T2U et U2T}

\begin{enumerate}[label=\alph*)]
    \item Appliquer $T2U_8(x)$ pour $x = -1$ (sur 8 bits). Vérifier avec la formule.
    \item Appliquer $T2U_8(x)$ pour $x = -128$.
    \item Appliquer $U2T_8(u)$ pour $u = 200$.
    \item Appliquer $U2T_8(u)$ pour $u = 127$.
\end{enumerate}

\textit{Rappel :} $T2U_w(x) = x + 2^w$ si $x < 0$, sinon $x$. \quad $U2T_w(u) = u - 2^w$ si $u > T_{max}$, sinon $u$.
\end{exercicebox}

\begin{exercicebox}{Exercice 10 : Pièges de comparaison}

Prédire le résultat (vrai=1 ou faux=0) de chaque expression C. Justifier.

\begin{lstlisting}[numbers=none]
int x = -1;
unsigned int y = 0;

a) x < y           // ?
b) x == 0xFFFFFFFF // ? (sur machine 32 bits)
c) -1 < 0U         // ?
d) 2147483647 > -2147483648        // ? (types int)
e) 2147483647U > -2147483648       // ? (unsigned à gauche)
\end{lstlisting}
\end{exercicebox}

%==============================================================================
\section{Extension de Signe}
%==============================================================================

\begin{exercicebox}{Exercice 11 : Extension 8 bits $\rightarrow$ 16 bits}

Effectuer l'extension de signe de 8 vers 16 bits pour :

\begin{enumerate}[label=\alph*)]
    \item \texttt{int8\_t x = -6} (soit $11111010_2$) $\rightarrow$ \texttt{int16\_t}
    \item \texttt{int8\_t x = 120} (soit $01111000_2$) $\rightarrow$ \texttt{int16\_t}
    \item \texttt{uint8\_t x = 250} (soit $11111010_2$) $\rightarrow$ \texttt{uint16\_t}
\end{enumerate}

Donner le résultat en binaire (16 bits) et en hexadécimal.
\end{exercicebox}

\begin{exercicebox}{Exercice 12 : Piège de l'extension avec cast}

Que valent les variables suivantes ? Expliquer chaque étape.

\begin{lstlisting}[numbers=none]
int8_t a = -1;                    // a = ?
uint32_t b = a;                   // b = ?
uint32_t c = (uint8_t)a;          // c = ?
\end{lstlisting}

\textit{Indice : dans le cas (c), le cast en \texttt{uint8\_t} est appliqué AVANT l'extension.}
\end{exercicebox}

%==============================================================================
\section{Dépassements (Overflow)}
%==============================================================================

\begin{exercicebox}{Exercice 13 : Overflow unsigned}

Calculer le résultat des opérations suivantes sur \texttt{uint8\_t} (8 bits, non signé). Utiliser la règle du modulo $2^8$.

\begin{enumerate}[label=\alph*)]
    \item $255 + 1$
    \item $250 + 10$
    \item $0 - 1$
    \item $100 - 150$
\end{enumerate}
\end{exercicebox}

\begin{exercicebox}{Exercice 14 : Overflow signed (Undefined Behavior)}

Pour chaque opération sur \texttt{int8\_t} (8 bits, signé), indiquer :
\begin{itemize}
    \item Si le résultat est \textbf{défini} ou \textbf{UB} (Undefined Behavior)
    \item Le résultat probable (si on suppose un wrap-around, même si c'est UB)
\end{itemize}

\begin{enumerate}[label=\alph*)]
    \item $127 + 1$
    \item $-128 - 1$
    \item $100 + 20$
    \item $-100 - 50$
\end{enumerate}
\end{exercicebox}

\begin{exercicebox}{Exercice 15 : Détection d'overflow AVANT l'opération}

Écrire les conditions qui permettent de détecter un overflow \textbf{avant} d'effectuer l'opération (pour éviter l'UB).

\begin{enumerate}[label=\alph*)]
    \item Pour l'addition \texttt{int32\_t a + b}, quand y a-t-il overflow positif ?
    \item Pour l'addition \texttt{int32\_t a + b}, quand y a-t-il overflow négatif ?
    \item Compléter le code :
\begin{lstlisting}[numbers=none]
int32_t a = ..., b = ...;
if (/* condition overflow positif */ || /* condition overflow négatif */) {
    printf("Overflow!\n");
} else {
    int32_t sum = a + b;  // Safe
}
\end{lstlisting}
\end{enumerate}
\end{exercicebox}

%==============================================================================
\section{IEEE 754 (Flottants)}
%==============================================================================

\begin{exercicebox}{Exercice 16 : Encodage IEEE 754 float (32 bits)}

Encoder les nombres suivants en IEEE 754 float (32 bits). Donner le résultat en hexadécimal.

\begin{enumerate}[label=\alph*)]
    \item $+1.0$
    \item $-1.0$
    \item $+12.5$
\end{enumerate}

Pour chaque nombre, détailler :
\begin{itemize}
    \item Signe $s$
    \item Conversion en binaire fractionnaire
    \item Normalisation ($1.xxx \times 2^n$)
    \item Exposant biaisé ($E + 127$)
    \item Fraction (23 bits, sans le "1" implicite)
\end{itemize}
\end{exercicebox}

\begin{exercicebox}{Exercice 17 : Décodage IEEE 754}

Le float \texttt{0xC1480000} est stocké en mémoire. Décoder ce nombre.

\begin{enumerate}[label=\alph*)]
    \item Écrire en binaire (32 bits).
    \item Identifier : signe (1 bit), exposant (8 bits), fraction (23 bits).
    \item Calculer l'exposant réel ($E = \text{exposant} - 127$).
    \item Reconstituer la mantisse $M = 1.\text{fraction}$.
    \item Calculer la valeur décimale : $V = (-1)^s \times M \times 2^E$.
\end{enumerate}
\end{exercicebox}

\begin{exercicebox}{Exercice 18 : Valeurs spéciales IEEE 754}

Compléter le tableau avec les représentations hexadécimales (32 bits) :

\begin{center}
\begin{tabular}{|l|c|c|c|c|}
\hline
\textbf{Valeur} & \textbf{s} & \textbf{exp (8 bits)} & \textbf{frac (23 bits)} & \textbf{Hex} \\
\hline
$+0$ & & & & \\
\hline
$-0$ & & & & \\
\hline
$+\infty$ & & & & \\
\hline
$-\infty$ & & & & \\
\hline
NaN (exemple) & & & & \\
\hline
\end{tabular}
\end{center}

\textit{Rappel : $\pm\infty$ a exposant = 255 et fraction = 0. NaN a exposant = 255 et fraction $\neq$ 0.}
\end{exercicebox}

\begin{exercicebox}{Exercice 19 : Pièges des flottants}

\begin{enumerate}[label=\alph*)]
    \item Pourquoi $0.1_{10}$ n'a-t-il pas de représentation exacte en binaire ?
    \item Que vaut \texttt{0.1f + 0.2f == 0.3f} en C ? Pourquoi ?
    \item Expliquer pourquoi \texttt{(3.14f + 1e10f) - 1e10f} ne donne pas \texttt{3.14f}.
    \item Quelle est la bonne façon de comparer deux flottants ?
\end{enumerate}
\end{exercicebox}

\begin{exercicebox}{Exercice 20 : Analyse de code}

Prédire la sortie de chaque \texttt{printf}. Justifier.

\begin{lstlisting}
#include <stdio.h>
#include <math.h>

int main() {
    float a = 0.1f + 0.2f;
    float b = 0.3f;
    printf("a == b ? %d\n", a == b);
    
    float x = 1.0f / 0.0f;
    printf("x = %f, isinf(x) = %d\n", x, isinf(x));
    
    float y = 0.0f / 0.0f;
    printf("y == y ? %d, isnan(y) = %d\n", y == y, isnan(y));
    
    return 0;
}
\end{lstlisting}
\end{exercicebox}

\end{document}
